% ============================================================
%  Probeklausur Template (my_dhbw Stil, klausur_*.tex)
%  Formell mit Deckblatt-Header; Lösungen als grüne tcolorbox.
%  Renderer: pdflatex (zweimal)
% ============================================================
\documentclass[11pt, a4paper]{article}

\usepackage[ngerman]{babel}
\usepackage{fontspec}
\setmainfont[
  Path=/usr/share/fonts/TTF/,
  Extension=.ttf,
  BoldFont=DejaVuSans-Bold,
  ItalicFont=DejaVuSans-Oblique,
  BoldItalicFont=DejaVuSans-BoldOblique
]{DejaVuSans}
\setsansfont[
  Path=/usr/share/fonts/TTF/,
  Extension=.ttf,
  BoldFont=DejaVuSans-Bold,
  ItalicFont=DejaVuSans-Oblique,
  BoldItalicFont=DejaVuSans-BoldOblique
]{DejaVuSans}
\setmonofont[Path=/usr/share/fonts/TTF/,Extension=.ttf]{DejaVuSansMono}
\usepackage[top=2cm, bottom=2cm, left=2.4cm, right=2.4cm, footskip=1cm]{geometry}
\usepackage{amsmath, amssymb}
\usepackage{xcolor}
\usepackage{enumitem}
\usepackage{fancyhdr}
\usepackage{lastpage}
\usepackage{tabularx}
\usepackage{booktabs}
\usepackage{microtype}
\usepackage{parskip}

\usepackage{array}
\usepackage{longtable}
\usepackage{ltablex}
\keepXColumns
\usepackage{colortbl}
\usepackage[most,breakable]{tcolorbox}
\usepackage{listings}
\usepackage[normalem]{ulem}
\usepackage{pdflscape}
\usepackage{titlesec}
\usepackage[colorlinks=true, linkcolor=accent, urlcolor=accent]{hyperref}

\renewcommand{\familydefault}{\sfdefault}

% ── Theme ─────────────────────────────────────────────────────
\definecolor{accent}{RGB}{0, 60, 113}
\definecolor{primaryblue}{RGB}{0, 60, 113}
\definecolor{loesung}{RGB}{0, 100, 0}
\definecolor{codebg}{RGB}{245, 245, 245}
\definecolor{figurebg}{RGB}{232, 241, 255}
\definecolor{tablehintbg}{RGB}{240, 255, 240}
\definecolor{solutionbg}{RGB}{240, 252, 240}
\definecolor{rulegray}{RGB}{180, 180, 180}
\definecolor{tableheadbg}{RGB}{0, 60, 113}

% ── Header/Footer ─────────────────────────────────────────────
\pagestyle{fancy}
\fancyhf{}
\renewcommand{\headrulewidth}{0.4pt}
\fancyhead[L]{\footnotesize\color{accent}Modulprüfung -- \nichttitle}
\fancyhead[R]{\footnotesize\color{accent}\nichtsubtitle}
\fancyfoot[L]{\footnotesize\color{accent}Matrikelnr.: \rule{3cm}{0.4pt}}
\fancyfoot[R]{\footnotesize\color{accent}Blatt \thepage\,/\,\pageref{LastPage}}

% ── Typografie ────────────────────────────────────────────────
\titleformat{\section}
    {\color{accent}\large\bfseries}{}{0pt}{}
    [\vspace{-4pt}\color{accent}\rule{\linewidth}{1.5pt}]
\titleformat{\subsection}
    {\normalsize\bfseries\color{accent!85!black}}{}{0pt}{}{}
\titleformat{\subsubsection}
    {\small\bfseries\color{accent!70!black}}{}{0pt}{}{}
\titlespacing*{\section}{0pt}{12pt}{4pt}
\titlespacing*{\subsection}{0pt}{8pt}{2pt}
\titlespacing*{\subsubsection}{0pt}{6pt}{1pt}

\setlength{\parindent}{0pt}
\setlength{\parskip}{6pt}
\renewcommand{\arraystretch}{1.25}
\setlist[itemize]{noitemsep, topsep=3pt, parsep=0pt, leftmargin=1.4em}
\setlist[enumerate]{noitemsep, topsep=3pt, parsep=0pt, leftmargin=1.6em}

% ── Boxen ─────────────────────────────────────────────────────
\newtcolorbox{figurebox}[1]{
  colback=figurebg, colframe=accent!50,
  title={Abbildung: #1}, fonttitle=\small\bfseries,
  breakable, before skip=6pt, after skip=6pt
}
\newtcolorbox{tablehintbox}[1]{
  colback=tablehintbg, colframe=green!50!black!50,
  title={Tabelle: #1}, fonttitle=\small\bfseries,
  breakable, before skip=6pt, after skip=6pt
}
\newtcolorbox{solutionbox}[1]{
  enhanced,
  colback=solutionbg, colframe=loesung,
  title={#1}, fonttitle=\small\bfseries,
  coltitle=white, colbacktitle=loesung,
  breakable, before skip=8pt, after skip=8pt
}

\lstset{
  basicstyle=\ttfamily\small,
  backgroundcolor=\color{codebg},
  breaklines=true,
  frame=single, framerule=0pt,
  xleftmargin=4pt, xrightmargin=4pt,
  aboveskip=6pt, belowskip=6pt,
  keepspaces=true
}

\providecommand{\nichttitle}{Probeklausur}
\providecommand{\nichtsubtitle}{}

\renewcommand{\nichttitle}{Relationen, Algebra, Diskrete Mathematik}
\renewcommand{\nichtsubtitle}{Probeklausur 3}


\begin{document}

% Deckblatt
\thispagestyle{empty}
\begin{center}
\vspace*{1cm}
{\Huge\bfseries\color{accent}Probeklausur}\\[8pt]
{\Large\color{accent}\nichttitle}\\[4pt]
\ifx\nichtsubtitle\empty\else
  {\large\color{accent!80!black}\nichtsubtitle}\\[6pt]
\fi
\color{accent}\rule{0.5\linewidth}{1pt}\color{black}
\end{center}

\vspace{1cm}
\begin{tabularx}{\textwidth}{@{}l X@{}}
\textbf{Studiengang:}      & \rule{6cm}{0.4pt} \\[6pt]
\textbf{Matrikelnummer:}    & \rule{6cm}{0.4pt} \\[6pt]
\textbf{Name, Vorname:}    & \rule{6cm}{0.4pt} \\[6pt]
\textbf{Datum:}             & \rule{6cm}{0.4pt} \\[6pt]
\textbf{Bearbeitungszeit:}  & 60 Minuten \\[6pt]
\textbf{Hilfsmittel:}       & Taschenrechner (nicht programmierbar) \\
\end{tabularx}

\vspace{2cm}
\begin{center}
{\large\itshape Viel Erfolg!}
\end{center}

\newpage

% ── BODY ──────────────────────────────────────────────────────

\section{Probeklausur 3 — Relationen, Algebra, Diskrete Mathematik}

\textbf{Bearbeitungszeit}: 60 Minuten | \textbf{Hilfsmittel}: keine | \textbf{Punkte}: 100



\noindent\textcolor{rulegray}{\rule{\linewidth}{0.4pt}}


\paragraph{Aufgabe 1 — Boolesche Algebra \textit{(18 Punkte)}}\mbox{}\\[2pt]

a) \textit{(3 P)} Notieren Sie beide Formen der DeMorgan-Gesetze sowie das Dualitätsprinzip in einem Satz.


b) \textit{(5 P)} Beweisen Sie das Absorptionsgesetz \textbf{a ⊗ (a ⊕ b) = a} ausschließlich aus den Axiomen (Kommutativ, Distributiv, Neutralelement, Komplement) sowie dem 0-1-Gesetz \textbf{0 ⊗ b = 0}. Geben Sie zu jeder Umformung das verwendete Gesetz an.


c) \textit{(10 P)} Gegeben sei die Boolesche Funktion f(a, b, c, d) durch folgende Wertetabelle (nur 1-Zeilen aufgelistet, alle übrigen Belegungen liefern 0):



{\small
\begin{tabularx}{\linewidth}{XXXXX}
\hline
\rowcolor{tableheadbg}
{\color{white}\textbf{a}} & {\color{white}\textbf{b}} & {\color{white}\textbf{c}} & {\color{white}\textbf{d}} & {\color{white}\textbf{f}} \\[2pt]
\hline
\endhead
\hline
\endfoot
0 & 0 & 0 & 0 & 1 \\
0 & 0 & 0 & 1 & 1 \\
1 & 0 & 0 & 0 & 1 \\
1 & 0 & 0 & 1 & 1 \\
1 & 0 & 1 & 1 & 1 \\
1 & 1 & 0 & 1 & 1 \\
1 & 1 & 1 & 1 & 1 \\
\end{tabularx}
}


\begin{itemize}
  \item (c1, 3 P) Geben Sie die kanonische disjunktive Normalform (kDN) an.
\begin{itemize}
  \item (c2, 5 P) Zeichnen Sie das KV-Diagramm (4 Variablen, Gray-Code) und identifizieren Sie alle maximalen Einserblöcke grafisch.
  \item (c3, 2 P) Lesen Sie aus dem KV-Diagramm die disjunktive Minimalform (DM) ab.
\end{itemize}
\end{itemize}


\noindent\textcolor{rulegray}{\rule{\linewidth}{0.4pt}}


\paragraph{Aufgabe 2 — Schaltnetze: 2-Bit-Addierer \textit{(14 Punkte)}}\mbox{}\\[2pt]

a) \textit{(4 P)} Stellen Sie die Wahrheitstabelle eines Volladdierers (Eingänge a, b, ü₁) vollständig auf und leiten Sie die Booleschen Ausdrücke für die Summe s und den Ausgangsübertrag ü₂ aus dem Lernzettel ab.


b) \textit{(6 P)} Konstruieren Sie ein Schaltnetz zur Addition zweier 2-Bit-Zahlen (a₁a₀) + (b₁b₀) → (s₂s₁s₀). Verwenden Sie genau \textbf{einen Halbaddierer (HA)} und \textbf{einen Volladdierer (VA)}. Skizzieren Sie das Blockschaltbild, beschriften Sie alle Signalleitungen (insbesondere den Übertragspfad) und geben Sie die resultierenden Booleschen Ausdrücke für s₀, s₁ und s₂ an.


c) \textit{(4 P)} Begründen Sie ausführlich, warum für die niederwertigste Stelle ein Halbaddierer ausreicht und kein Volladdierer benötigt wird. Erläutern Sie dabei, was an dieser Stelle bezüglich des Eingangsübertrags gilt.



\noindent\textcolor{rulegray}{\rule{\linewidth}{0.4pt}}


\paragraph{Aufgabe 3 — Graphen: Eigenschaften und Darstellung \textit{(12 Punkte)}}\mbox{}\\[2pt]

Gegeben sei der ungerichtete Graph G = (V, E) mit V = \{A, B, C, D, E\} und folgenden gewichteten Kanten:


E = \{ (A,B,3), (A,C,2), (B,C,4), (B,D,1), (C,E,5), (D,E,2) \}


a) \textit{(4 P)} Erstellen Sie sowohl die \textbf{Adjazenzmatrix} (5 × 5, mit Gewichten) als auch die \textbf{Adjazenzliste} des Graphen.


b) \textit{(4 P)} Bestimmen Sie für jeden Knoten den Knotengrad. Untersuchen Sie anschließend, ob G ein Multigraph, ein vollständiger Graph oder ein bipartiter Graph ist. Begründen Sie jede Antwort.


c) \textit{(4 P)} Existiert in G ein \textbf{Eulerscher Pfad}? Beantworten Sie die Frage mithilfe der Knotengrade aus (b), benennen Sie das verwendete Kriterium und geben Sie ggf. Anfangs- und Endknoten eines möglichen Eulerschen Pfades an.



\noindent\textcolor{rulegray}{\rule{\linewidth}{0.4pt}}


\paragraph{Aufgabe 4 — Dijkstra-Algorithmus \textit{(20 Punkte)}}\mbox{}\\[2pt]

Gegeben sei der ungerichtete, gewichtete Graph mit V = \{1, 2, 3, 4, 5\} und Kanten:


(1,2,4), (1,3,2), (2,3,1), (2,4,5), (3,4,8), (3,5,10), (4,5,2)


a) \textit{(3 P)} Geben Sie die Initialwerte für d(v), π(v) und α(v) für alle v ∈ V bei Startknoten s = 1 an.


b) \textit{(12 P)} Führen Sie den Dijkstra-Algorithmus vollständig durch. Dokumentieren Sie nach jeder Knotenwahl in einer Tabelle die aktuellen Werte (d, α) für alle Knoten sowie den als nächstes zu besuchenden Knoten.


c) \textit{(3 P)} Rekonstruieren Sie den kürzesten Pfad 1 → 5 anhand der α-Werte.


d) \textit{(2 P)} Angenommen, die Kante (3, 5) hätte das Gewicht \textbf{−3} statt 10. Welcher Schritt im Dijkstra-Algorithmus würde versagen? Welcher Algorithmus aus dem Vergleichsabschnitt wäre stattdessen geeignet?



\noindent\textcolor{rulegray}{\rule{\linewidth}{0.4pt}}


\paragraph{Aufgabe 5 — A\textit{-Algorithmus und Heuristik-Bewertung }(14 Punkte)*}\mbox{}\\[2pt]

a) \textit{(3 P)} Erläutern Sie den Hauptunterschied zwischen Dijkstra und A* in Bezug auf (i) die zusätzlich gespeicherte Knotengröße und (ii) das Kriterium der Knotenwahl in der Hauptschleife.


b) \textit{(4 P)} Was ist eine Heuristik h(v, z) im Kontext von A\textit{? Welche Eigenschaft muss sie erfüllen, damit A} garantiert die optimale Lösung findet? Was geschieht laut Lernzettel, wenn diese Eigenschaft verletzt wird?


c) \textit{(7 P)} Für den Graphen aus Aufgabe 4 (Ziel z = 5) sei folgende Heuristik gegeben:



{\small
\begin{tabularx}{\linewidth}{XX}
\hline
\rowcolor{tableheadbg}
{\color{white}\textbf{Knoten v}} & {\color{white}\textbf{h(v, 5)}} \\[2pt]
\hline
\endhead
\hline
\endfoot
1 & 9 \\
2 & 6 \\
3 & 8 \\
4 & 1 \\
5 & 0 \\
\end{tabularx}
}


Bestimmen Sie für jeden Knoten v die \textbf{tatsächliche} kürzeste Distanz d\textit{(v, 5) und entscheiden Sie tabellarisch, ob h(v, 5) ≤ d}(v, 5) gilt. Begründen Sie abschließend, ob A* mit dieser Heuristik die optimale Lösung liefern wird.



\noindent\textcolor{rulegray}{\rule{\linewidth}{0.4pt}}


\paragraph{Aufgabe 6 — TSP: Nearest-Neighbour \& 2-Opt \textit{(22 Punkte)}}\mbox{}\\[2pt]

Gegeben seien fünf Städte A, B, C, D, E mit folgender symmetrischer Distanzmatrix:



{\small
\begin{tabularx}{\linewidth}{XXXXXX}
\hline
\rowcolor{tableheadbg}
{\color{white}\textbf{}} & {\color{white}\textbf{A}} & {\color{white}\textbf{B}} & {\color{white}\textbf{C}} & {\color{white}\textbf{D}} & {\color{white}\textbf{E}} \\[2pt]
\hline
\endhead
\hline
\endfoot
A & 0 & 2 & 4 & 1 & 7 \\
B & 2 & 0 & 3 & 5 & 4 \\
C & 4 & 3 & 0 & 6 & 2 \\
D & 1 & 5 & 6 & 0 & 8 \\
E & 7 & 4 & 2 & 8 & 0 \\
\end{tabularx}
}


a) \textit{(4 P)} Definieren Sie das TSP formal (Eingabe, gesuchte Ausgabe, Zielgröße). Nennen Sie zwei aus dem Lernzettel bekannte Anwendungsgebiete und ein verwandtes Problem.


b) \textit{(8 P)} Wenden Sie die \textbf{Nearest-Neighbour-Heuristik (NNA)} mit Startknoten \textbf{A} an. Dokumentieren Sie den aktuellen Knoten, die zur Auswahl stehenden Kanten und die jeweils gewählte Kante in jedem Schritt. Geben Sie am Ende die vollständige Tour und ihre Gesamtlänge an.


c) \textit{(8 P)} Wenden Sie auf das NNA-Ergebnis aus (b) \textbf{einen 2-Opt-Schritt} an: Prüfen Sie alle Inversionen einer zusammenhängenden Subsequenz der \textbf{Länge 2} (also Vertauschungen zweier aufeinanderfolgender innerer Knoten). Berechnen Sie für jede Variante die neue Tourlänge und geben Sie die beste Verbesserung samt resultierender Tour an.


d) \textit{(2 P)} Begründen Sie kurz, warum NNA das globale Optimum nicht garantieren kann, und nennen Sie eine im Lernzettel genannte einfache Erweiterung, die die Lösungsqualität verbessert.



\noindent\textcolor{rulegray}{\rule{\linewidth}{0.4pt}}



\section{Musterlösung Klausur 3}

\paragraph{Aufgabe 1 \textit{(18 Punkte)}}\mbox{}\\[2pt]

\textbf{a) DeMorgan und Dualität (3 P)}


\begin{itemize}
  \item ∼(a ⊗ b) = ∼a ⊕ ∼b \textit{(1 P)}
  \item ∼(a ⊕ b) = ∼a ⊗ ∼b \textit{(1 P)}
  \item Dualitätsprinzip: Vertauscht man in einer gültigen Aussage gleichzeitig ⊗ ↔ ⊕ und 0 ↔ 1, so entsteht erneut eine gültige Aussage. \textit{(1 P)}
\end{itemize}

\textbf{b) Beweis Absorption a ⊗ (a ⊕ b) = a (5 P)}



{\small
\begin{tabularx}{\linewidth}{XXX}
\hline
\rowcolor{tableheadbg}
{\color{white}\textbf{Schritt}} & {\color{white}\textbf{Ausdruck}} & {\color{white}\textbf{Begründung}} \\[2pt]
\hline
\endhead
\hline
\endfoot
1 & a ⊗ (a ⊕ b) & Ausgangsausdruck \\
2 & (a ⊕ 0) ⊗ (a ⊕ b) & Neutralelement: a = a ⊕ 0 \\
3 & a ⊕ (0 ⊗ b) & Distributivgesetz Form 2 rückwärts \\
4 & a ⊕ 0 & 0-1-Gesetz: 0 ⊗ b = 0 \\
5 & a & Neutralelement: a ⊕ 0 = a \\
\end{tabularx}
}


\textit{Bewertung: pro korrekt benannter Umformung 1 P (5 × 1 P).} Alternative Beweise (z. B. via Distributiv Form 1 + Idempotenz + Absorption Form 2) werden nicht akzeptiert, sofern sie ein noch nicht bewiesenes Absorptionsgesetz benutzen (zirkulär).


\textbf{c) Boolesche Funktion f(a, b, c, d)}


\textbf{c1) kDN (3 P)} — pro Zeile mit f = 1 ein Minterm; gesamt 7 Minterme:


f = a′b′c′d′ ∨ a′b′c′d ∨ ab′c′d′ ∨ ab′c′d ∨ ab′cd ∨ abc′d ∨ abcd


\textit{Bewertung: 7 korrekte Minterme = 3 P; pro fehlendem/falschem Minterm −0,5 P.}


\textbf{c2) KV-Diagramm (5 P)} — Spalten cd ∈ \{00, 01, 11, 10\}, Zeilen ab ∈ \{00, 01, 11, 10\}:



{\small
\begin{tabularx}{\linewidth}{XXXXX}
\hline
\rowcolor{tableheadbg}
{\color{white}\textbf{ab\textbackslash{}cd}} & {\color{white}\textbf{00}} & {\color{white}\textbf{01}} & {\color{white}\textbf{11}} & {\color{white}\textbf{10}} \\[2pt]
\hline
\endhead
\hline
\endfoot
\textbf{00} & 1 & 1 & 0 & 0 \\
\textbf{01} & 0 & 0 & 0 & 0 \\
\textbf{11} & 0 & 1 & 1 & 0 \\
\textbf{10} & 1 & 1 & 1 & 0 \\
\end{tabularx}
}


Maximale Einserblöcke:

\begin{itemize}
  \item \textbf{Block 1} (4er-Block): Zeilen ab ∈ \{00, 10\}, Spalten cd ∈ \{00, 01\} → b und c konstant 0; a und d wechseln → bleibt \textbf{b′c′}.
  \item \textbf{Block 2} (4er-Block): Zeilen ab ∈ \{11, 10\}, Spalten cd ∈ \{01, 11\} → a konstant 1, d konstant 1; b und c wechseln → bleibt \textbf{ad}.
\end{itemize}

Beide Blöcke sind notwendig; jede 1-Zelle wird abgedeckt. \textit{(KV korrekt 2 P, beide Blöcke korrekt identifiziert je 1,5 P)}


\textbf{c3) Disjunktive Minimalform (2 P)}


f = b′c′ ∨ ad


\textit{Bewertung: vollständig korrekt 2 P; fehlender Term oder falsche Variable −1 P.}



\noindent\textcolor{rulegray}{\rule{\linewidth}{0.4pt}}


\paragraph{Aufgabe 2 \textit{(14 Punkte)}}\mbox{}\\[2pt]

\textbf{a) Volladdierer-Wahrheitstabelle und Funktionen (4 P)}



{\small
\begin{tabularx}{\linewidth}{XXXXX}
\hline
\rowcolor{tableheadbg}
{\color{white}\textbf{a}} & {\color{white}\textbf{b}} & {\color{white}\textbf{ü₁}} & {\color{white}\textbf{s}} & {\color{white}\textbf{ü₂}} \\[2pt]
\hline
\endhead
\hline
\endfoot
0 & 0 & 0 & 0 & 0 \\
0 & 0 & 1 & 1 & 0 \\
0 & 1 & 0 & 1 & 0 \\
0 & 1 & 1 & 0 & 1 \\
1 & 0 & 0 & 1 & 0 \\
1 & 0 & 1 & 0 & 1 \\
1 & 1 & 0 & 0 & 1 \\
1 & 1 & 1 & 1 & 1 \\
\end{tabularx}
}


Funktionen:

\begin{itemize}
  \item s = a ⊕ b ⊕ ü₁
  \item ü₂ = ab ∨ aü₁ ∨ bü₁
\end{itemize}

\textit{Bewertung: Tabelle korrekt 2 P; s 1 P; ü₂ 1 P.}


\textbf{b) 2-Bit-Addierer (a₁a₀) + (b₁b₀) (6 P)}


Blockschaltbild:


\begin{lstlisting}
   a₀ ──┐         a₁ ──┐
        │              │
   b₀ ──┤    HA   b₁ ──┤    VA   ── s₂ (= ü₂)
        │ ──┐          │
        ▼   │ s₀       ▼
        ü ──┘─────────→ü₁                ── s₁
\end{lstlisting}


\begin{itemize}
  \item HA für die niederwertigste Stelle: Eingänge a₀, b₀ → Ausgänge s₀, ü.
  \item ü wird als Eingangsübertrag ü₁ in den VA geführt.
  \item VA für die obere Stelle: Eingänge a₁, b₁, ü₁ → Ausgänge s₁, ü₂; ü₂ wird zu s₂.
\end{itemize}

Boolesche Ausdrücke:

\begin{itemize}
  \item s₀ = a₀ ⊕ b₀
  \item ü = a₀ ⊗ b₀
  \item s₁ = a₁ ⊕ b₁ ⊕ ü
  \item s₂ = a₁b₁ ∨ a₁ü ∨ b₁ü
\end{itemize}

\textit{Bewertung: Skizze mit korrekter Übertragsleitung 2 P; korrekte Zuordnung HA/VA 1 P; vier Boolesche Ausdrücke je 0,75 P (s₀, ü/Übertrag, s₁, s₂) = 3 P.}


\textbf{c) Begründung HA für die LSB (4 P)}


An der niederwertigsten Stelle gibt es \textbf{keinen Eingangsübertrag} aus einer noch niedrigeren Stelle (die Addition beginnt dort). Ein Volladdierer hätte in seinem dritten Eingang ü₁ stets den konstanten Wert 0. Mit ü₁ = 0 vereinfacht sich die VA-Funktion zu:

\begin{itemize}
  \item s = a ⊕ b ⊕ 0 = a ⊕ b
  \item ü₂ = ab ∨ a·0 ∨ b·0 = ab
\end{itemize}

Das entspricht genau der Funktion eines Halbaddierers. Daher ist der HA an der LSB ausreichend, ohne Funktionalität zu verlieren — er spart Gatter und damit Fläche/Verbrauch ein.


\textit{Bewertung: kein Eingangsübertrag an LSB 1 P; Reduktion VA → HA gezeigt 2 P; Effizienzargument 1 P.}



\noindent\textcolor{rulegray}{\rule{\linewidth}{0.4pt}}


\paragraph{Aufgabe 3 \textit{(12 Punkte)}}\mbox{}\\[2pt]

\textbf{a) Darstellungen (4 P)}


Adjazenzmatrix (Reihenfolge A, B, C, D, E; Gewichte; 0 = keine Kante):



{\small
\begin{tabularx}{\linewidth}{XXXXXX}
\hline
\rowcolor{tableheadbg}
{\color{white}\textbf{}} & {\color{white}\textbf{A}} & {\color{white}\textbf{B}} & {\color{white}\textbf{C}} & {\color{white}\textbf{D}} & {\color{white}\textbf{E}} \\[2pt]
\hline
\endhead
\hline
\endfoot
\textbf{A} & 0 & 3 & 2 & 0 & 0 \\
\textbf{B} & 3 & 0 & 4 & 1 & 0 \\
\textbf{C} & 2 & 4 & 0 & 0 & 5 \\
\textbf{D} & 0 & 1 & 0 & 0 & 2 \\
\textbf{E} & 0 & 0 & 5 & 2 & 0 \\
\end{tabularx}
}


Adjazenzliste:

\begin{itemize}
  \item A → [(B, 3), (C, 2)]
  \item B → [(A, 3), (C, 4), (D, 1)]
  \item C → [(A, 2), (B, 4), (E, 5)]
  \item D → [(B, 1), (E, 2)]
  \item E → [(C, 5), (D, 2)]
\end{itemize}

\textit{Bewertung: Matrix symmetrisch und korrekt 2 P; Adjazenzliste korrekt 2 P.}


\textbf{b) Knotengrade \& Eigenschaften (4 P)}


Knotengrade: deg(A) = 2, deg(B) = 3, deg(C) = 3, deg(D) = 2, deg(E) = 2 \textit{(1 P)}


\begin{itemize}
  \item \textbf{Multigraph?} Nein. Zwischen je zwei Knoten existiert höchstens eine Kante. \textit{(1 P)}
  \item \textbf{Vollständig?} Nein. Ein vollständiger Graph mit 5 Knoten besäße C(5,2) = 10 Kanten; hier sind es nur 6. \textit{(1 P)}
  \item \textbf{Bipartit?} Nein. Die Knoten \{A, B, C\} bilden ein Dreieck (A–B, B–C, A–C) — ein ungerader Zyklus. Ein bipartiter Graph enthält keine ungeraden Zyklen. \textit{(1 P)}
\end{itemize}

\textbf{c) Eulerscher Pfad (4 P)}


Kriterium: Ein zusammenhängender ungerichteter Graph besitzt genau dann einen Eulerschen Pfad, wenn die Anzahl der Knoten mit ungeradem Grad \textbf{0 oder 2} beträgt. \textit{(2 P)}


Knoten ungeraden Grades: B (Grad 3) und C (Grad 3) → genau 2 → \textbf{Eulerscher Pfad existiert}. \textit{(1 P)}


Anfangs- und Endknoten: B und C (in beliebiger Richtung). \textit{(1 P)}



\noindent\textcolor{rulegray}{\rule{\linewidth}{0.4pt}}


\paragraph{Aufgabe 4 \textit{(20 Punkte)}}\mbox{}\\[2pt]

\textbf{a) Initialisierung (3 P)}



{\small
\begin{tabularx}{\linewidth}{XXXX}
\hline
\rowcolor{tableheadbg}
{\color{white}\textbf{v}} & {\color{white}\textbf{d(v)}} & {\color{white}\textbf{π(v)}} & {\color{white}\textbf{α(v)}} \\[2pt]
\hline
\endhead
\hline
\endfoot
1 & 0 & unvisited & none \\
2 & ∞ & unvisited & none \\
3 & ∞ & unvisited & none \\
4 & ∞ & unvisited & none \\
5 & ∞ & unvisited & none \\
\end{tabularx}
}


Aktueller Knoten: u = 1.


\textit{Bewertung: korrekte d-Werte 1 P; π = unvisited für alle 1 P; s.d = 0 / s als Start 1 P.}


\textbf{b) Schrittweise Ausführung (12 P)}


Notation: nach jeder Visite wird der besuchte Knoten markiert; Tabellenzeile zeigt Werte \textbf{nach} Update der Nachbarn.


\textbf{Schritt 1 — Besuche Knoten 1 (d = 0):}

Updates: 2.d = min(∞, 0+4) = 4 (α=1), 3.d = min(∞, 0+2) = 2 (α=1).



{\small
\begin{tabularx}{\linewidth}{XXXXXX}
\hline
\rowcolor{tableheadbg}
{\color{white}\textbf{v}} & {\color{white}\textbf{1}} & {\color{white}\textbf{2}} & {\color{white}\textbf{3}} & {\color{white}\textbf{4}} & {\color{white}\textbf{5}} \\[2pt]
\hline
\endhead
\hline
\endfoot
d & 0 & 4 & 2 & ∞ & ∞ \\
α & – & 1 & 1 & – & – \\
\end{tabularx}
}


Nächster: \textbf{3} (kleinstes d unter unvisited = 2). \textit{(2,5 P)}


\textbf{Schritt 2 — Besuche Knoten 3 (d = 2):}

Updates Nachbarn (1 visited): 2.d = min(4, 2+1) = 3 (α=3), 4.d = min(∞, 2+8) = 10 (α=3), 5.d = min(∞, 2+10) = 12 (α=3).



{\small
\begin{tabularx}{\linewidth}{XXXXXX}
\hline
\rowcolor{tableheadbg}
{\color{white}\textbf{v}} & {\color{white}\textbf{1}} & {\color{white}\textbf{2}} & {\color{white}\textbf{3}} & {\color{white}\textbf{4}} & {\color{white}\textbf{5}} \\[2pt]
\hline
\endhead
\hline
\endfoot
d & 0 & 3 & 2 & 10 & 12 \\
α & – & 3 & 1 & 3 & 3 \\
\end{tabularx}
}


Nächster: \textbf{2} (d = 3). \textit{(2,5 P)}


\textbf{Schritt 3 — Besuche Knoten 2 (d = 3):}

Updates Nachbarn (1, 3 visited): 4.d = min(10, 3+5) = 8 (α=2).



{\small
\begin{tabularx}{\linewidth}{XXXXXX}
\hline
\rowcolor{tableheadbg}
{\color{white}\textbf{v}} & {\color{white}\textbf{1}} & {\color{white}\textbf{2}} & {\color{white}\textbf{3}} & {\color{white}\textbf{4}} & {\color{white}\textbf{5}} \\[2pt]
\hline
\endhead
\hline
\endfoot
d & 0 & 3 & 2 & 8 & 12 \\
α & – & 3 & 1 & 2 & 3 \\
\end{tabularx}
}


Nächster: \textbf{4} (d = 8). \textit{(2,5 P)}


\textbf{Schritt 4 — Besuche Knoten 4 (d = 8):}

Updates Nachbarn (2, 3 visited): 5.d = min(12, 8+2) = 10 (α=4).



{\small
\begin{tabularx}{\linewidth}{XXXXXX}
\hline
\rowcolor{tableheadbg}
{\color{white}\textbf{v}} & {\color{white}\textbf{1}} & {\color{white}\textbf{2}} & {\color{white}\textbf{3}} & {\color{white}\textbf{4}} & {\color{white}\textbf{5}} \\[2pt]
\hline
\endhead
\hline
\endfoot
d & 0 & 3 & 2 & 8 & 10 \\
α & – & 3 & 1 & 2 & 4 \\
\end{tabularx}
}


Nächster: \textbf{5} (d = 10). \textit{(2,5 P)}


\textbf{Schritt 5 — Besuche Knoten 5 (d = 10):} Keine unvisited Nachbarn mehr → fertig. \textit{(2 P)}


\textbf{c) Pfadrekonstruktion 1 → 5 (3 P)}


α-Kette: 5 ← 4 ← 2 ← 3 ← 1. Kürzester Pfad: \textbf{1 → 3 → 2 → 4 → 5} mit Länge 2 + 1 + 5 + 2 = 10. \textit{(Pfad korrekt 2 P; Länge 1 P)}


\textbf{d) Negative Kantengewichte (2 P)}


Mit w(3,5) = −3 würde die Update-Regel verletzt: Knoten 3 wird in Schritt 2 als visited markiert; ein späterer kürzerer Weg über die negative Kante könnte 3 nicht mehr aktualisieren, weil Dijkstra besuchte Knoten nicht mehr re-prüft. Der Algorithmus liefert dann i. A. nicht die optimale Distanz. Geeignet ist \textbf{Bellman–Ford} (laut Vergleichstabelle der einzige Single-Source-Algorithmus, der negative Gewichte erlaubt). \textit{(Begründung 1 P; Bellman–Ford 1 P)}



\noindent\textcolor{rulegray}{\rule{\linewidth}{0.4pt}}


\paragraph{Aufgabe 5 \textit{(14 Punkte)}}\mbox{}\\[2pt]

\textbf{a) Unterschied Dijkstra ↔ A* (3 P)}


(i) A\textit{ speichert pro Knoten v zusätzlich zur Distanz d(v) eine }\textit{Priorität}\textit{ p(v) = d(v) + h(v, z). }(1,5 P)*

(ii) In der Hauptschleife wählt A\textit{ als nächsten aktuellen Knoten denjenigen unvisited Knoten mit }\textit{minimaler Priorität p}\textit{ (statt minimalem d wie bei Dijkstra). }(1,5 P)*


\textbf{b) Heuristik \& Zulässigkeit (4 P)}


Eine Heuristik h(v, z) ist eine \textbf{Schätzung der Restdistanz} vom Knoten v zum Zielknoten z (z. B. Luftlinie). \textit{(1,5 P)}


Damit A\textit{ das Optimum garantiert findet, muss h }\textit{zulässig (admissible)}\textit{ sein: h(v, z) darf die tatsächliche kürzeste Distanz v → z }\textit{nie überschätzen}\textit{, d. h. h(v, z) ≤ d}(v, z). \textit{(1,5 P)}


Wird die Eigenschaft verletzt (Heuristik überschätzt), liefert A\textit{ laut Lernzettel u. U. eine }\textit{suboptimale Lösung}\textit{, weil ein an sich optimaler Knoten zu spät bzw. gar nicht expandiert wird. }(1 P)*


\textbf{c) Zulässigkeitsprüfung (7 P)}


Tatsächliche kürzeste Distanzen v → 5 anhand des Graphen aus Aufgabe 4:

\begin{itemize}
  \item d*(1, 5) = 10 (Pfad 1-3-2-4-5)
  \item d*(2, 5) = 7 (Pfad 2-4-5)
  \item d*(3, 5) = 8 (Pfad 3-2-4-5: 1+5+2)
  \item d*(4, 5) = 2 (Pfad 4-5)
  \item d*(5, 5) = 0
\end{itemize}


{\small
\begin{tabularx}{\linewidth}{XXXX}
\hline
\rowcolor{tableheadbg}
{\color{white}\textbf{v}} & {\color{white}\textbf{h(v, 5)}} & {\color{white}\textbf{d*(v, 5)}} & {\color{white}\textbf{h ≤ d* ?}} \\[2pt]
\hline
\endhead
\hline
\endfoot
1 & 9 & 10 & ✓ \\
2 & 6 & 7 & ✓ \\
3 & 8 & 8 & ✓ (=) \\
4 & 1 & 2 & ✓ \\
5 & 0 & 0 & ✓ \\
\end{tabularx}
}


Alle Heuristikwerte erfüllen h(v, 5) ≤ d\textit{(v, 5) → die Heuristik ist }\textit{zulässig}\textit{. Folglich wird A} mit dieser Heuristik den optimalen Pfad 1 → 5 mit Gesamtkosten 10 finden.


\textit{Bewertung: korrekte d}-Werte je 1 P (5 × 1 P), korrekte Tabellenauswertung 1 P, korrekte Schlussfolgerung 1 P.*



\noindent\textcolor{rulegray}{\rule{\linewidth}{0.4pt}}


\paragraph{Aufgabe 6 \textit{(22 Punkte)}}\mbox{}\\[2pt]

\textbf{a) TSP-Definition (4 P)}


\textbf{Eingabe}: Ein vollständiger gewichteter Graph G = (V, E, W) mit |V| = n Knoten (Städte) und nicht-negativen Kantengewichten (Distanzen). \textit{(1 P)}

\textbf{Gesuchte Ausgabe}: Eine geschlossene Tour, die \textbf{jede Stadt genau einmal} besucht und zur Startstadt zurückkehrt (Hamiltonscher Zyklus). \textit{(1 P)}

\textbf{Zielgröße}: Minimierung der Summe der Kantengewichte entlang der Tour. \textit{(0,5 P)}


Anwendungen: Logistik / Vehicle-Routing-Problem; Mikrochip-Design / Leiterplattenbohrung; Genom-Sequenzierung \textit{(2 davon, je 0,5 P = 1 P)}.

Verwandtes Problem: TSP ohne Rückkehr / VRP / Traveling Thief / stochastisches TSP \textit{(0,5 P)}.


\textbf{b) Nearest-Neighbour-Heuristik (NNA) ab A (8 P)}



{\footnotesize
\begin{tabularx}{\linewidth}{XXXXXXX}
\hline
\rowcolor{tableheadbg}
{\color{white}\textbf{Schritt}} & {\color{white}\textbf{Aktuell}} & {\color{white}\textbf{Unvisited}} & {\color{white}\textbf{Distanzen (Auswahl)}} & {\color{white}\textbf{Gewählt}} & {\color{white}\textbf{Tour bisher}} & {\color{white}\textbf{Länge}} \\[2pt]
\hline
\endhead
\hline
\endfoot
1 & A & \{B, C, D, E\} & B:2, C:4, \textbf{D:1}, E:7 & A → D & A–D & 1 \\
2 & D & \{B, C, E\} & B:5, C:6, E:8 & D → B (5) & A–D–B & 6 \\
3 & B & \{C, E\} & C:3, E:4 & B → C (3) & A–D–B–C & 9 \\
4 & C & \{E\} & E:2 & C → E (2) & A–D–B–C–E & 11 \\
5 & E & – (Rückkehr A) & A:7 & E → A (7) & A–D–B–C–E–A & \textbf{18} \\
\end{tabularx}
}


\textbf{Tour: A → D → B → C → E → A}, Gesamtlänge \textbf{1 + 5 + 3 + 2 + 7 = 18}.


\textit{Bewertung: pro Schritt mit korrekter Auswahl und Begründung 1 P (5 × 1 P); korrekte Endtour 1 P; korrekte Länge 2 P.}


\textbf{c) 2-Opt-Schritt (8 P)}


Innere Knoten der NNA-Tour: D, B, C, E. Inversionen einer Subsequenz der Länge 2:



{\small
\begin{tabularx}{\linewidth}{XXXXX}
\hline
\rowcolor{tableheadbg}
{\color{white}\textbf{Inversion}} & {\color{white}\textbf{Neue Tour}} & {\color{white}\textbf{Längenrechnung}} & {\color{white}\textbf{Länge}} & {\color{white}\textbf{Δ vs. 18}} \\[2pt]
\hline
\endhead
\hline
\endfoot
(D, B) & A–B–D–C–E–A & 2+5+6+2+7 & 22 & +4 \\
(B, C) & A–D–C–B–E–A & 1+6+3+4+7 & 21 & +3 \\
(C, E) & A–D–B–E–C–A & 1+5+4+2+4 & \textbf{16} & \textbf{−2} \\
\end{tabularx}
}


Beste Verbesserung: Inversion \textbf{(C, E)} liefert die neue Tour \textbf{A → D → B → E → C → A} mit Länge \textbf{16} (Verbesserung um 2). \textit{(je Berechnung 2 P; korrekte Auswahl + neue Tour 2 P)}


\textbf{d) Suboptimalität von NNA (2 P)}


NNA wählt in jedem Schritt nur die lokal kürzeste Kante (Greedy / Best-First) und berücksichtigt keine späteren Kosten. Insbesondere kann die letzte Rückkehrkante zur Startstadt sehr lang werden (im Beispiel E→A mit 7), was eine vermeintlich kostengünstige Auswahl im Nachhinein teuer macht. Außerdem ist das Ergebnis \textbf{startknotenabhängig}.


Verbesserung laut Lernzettel: \textbf{Jeden Knoten als Startknoten ausprobieren} und die kürzeste der n NNA-Touren wählen (Laufzeit × n, weiterhin keine Optimalitätsgarantie).


\textit{Bewertung: Begründung Suboptimalität 1 P; Erweiterung 1 P.}





% ──────────────────────────────────────────────────────────────

\end{document}
